\section{Grouping cap conditions}
\label{sec:blocks}

We next group the cap conditions into blocks and choose a measurable set
inside each common intersection.  A point in this set meets every cap
condition in its block.  We need a lower bound on the product of the chosen
measures and control of their mutual overlaps.  Together these bounds will
control the self-overlap of the Poisson witnesses in
Section~\ref{sec:poisson}.  The merger construction and the weighted
row estimate follow Li--Liu
\cite[Proposition~9.1 and Lemmas~9.4--9.5]{LiLiu2026Lattice}.
We give the argument for arbitrary indexed lists, retaining repeated members.

\begin{proposition}[Block estimates]\label{prop:blocks}
Fix constants $C_0\ge1$, $c_0>0$, and $L\ge1$.
There are constants $B_0,K,C\ge1$ and $n_1$, depending only on these
three constants, with the following property.
Let $n\ge n_1$, let $(X,\mu)$ be an atomless finite measure space, and let
$(C_i)_{i\in I}$ be an indexed list of measurable sets with
$0<c_i:=\mu(C_i)<\infty$ and $1\le m:=|I|\le n^2$.
Set
\[
 \bar c_i=\min(c_i,1/2),\qquad s_i=-\log\bar c_i,
\]
and let $d$ be a pseudometric on $I$.  For nonempty $S\subseteq I$, write
\[
 r(S)=\min_{j\in S}\max_{i\in S}d(i,j)^2,
\]
and choose an index $i(S)\in S$ minimizing the uncapped measure $c_i$.
Assume that
\begin{equation}\label{eq:block-affinity-input}
 \frac{\mu(C_i\cap C_j)}{\sqrt{\bar c_i\bar c_j}}
 \le C_0e^{-c_0d(i,j)^2}\qquad(i,j\in I),
\end{equation}
and that, whenever $r(S)\le B_0\log n$,
\begin{equation}\label{eq:block-lens-input}
 \mu\left(\bigcap_{i\in S}C_i\right)
 \ge e^{-L(1+r(S)+\log|S|)}\bar c_{i(S)}.
\end{equation}
Then there is a partition $\mathcal P$ of $I$ and measurable sets
$E_S\subseteq\bigcap_{i\in S}C_i$, $S\in\mathcal P$, of measures
$\lambda_S=\mu(E_S)$ such that
\begin{equation}\label{eq:block-anchor-loss}
 n^{-K}\bar c_{i(S)}\le\lambda_S\le\frac12.
\end{equation}
Writing $q_S=-\log\lambda_S$ and
\[
 w_{ST}=
 \frac{\mu(E_S\cap E_T)}
 {\sqrt{\lambda_S(1+\lambda_S)\lambda_T(1+\lambda_T)}},
\]
we have
\begin{align}
 \sum_{S\in\mathcal P}q_S
 &\le\sum_{i\in I}s_i+\frac{m}{10},
 \label{eq:block-log-cost}\\
 \sum_{S\in\mathcal P}\log(1+\lambda_S^{-1})
 &\le C\left(m+\sum_{i\in I}s_i\right),
 \label{eq:block-first-budget}\\
 \sum_{S\in\mathcal P}
 \log\left(1+\sum_{T\in\mathcal P\setminus\{S\}}w_{ST}\right)
 &\le C\left(m+\sum_{i\in I}s_i\right).
 \label{eq:block-row-budget}
\end{align}
\end{proposition}

\begin{proof}
Set $\zeta=1/10$ and $R_*=B_0\log n$; we will choose $B_0$ after
the constants governing the row estimate.  Define
\[
 \ell(S)=
 \begin{cases}
 0,&|S|=1,\\
 L(1+r(S)+\log|S|),&|S|\ge2.
 \end{cases}
\]
For each final block $S$, we will choose a common subset of intensity
$e^{-\ell(S)}\bar c_{i(S)}$.  Thus $\ell(S)$ is the logarithmic loss
assigned to that block.  The merger condition below keeps the total loss
at most $\zeta m$, while terminality will control the number of nearby
blocks with small loss.

\emph{Constructing the partition.}
Start with the singleton partition.  Whenever a subfamily of at least two
current blocks has union $W$ satisfying
\begin{equation}\label{eq:block-merge}
 r(W)\le R_*,\qquad \ell(W)\le\zeta|W|,
\end{equation}
replace that entire subfamily by $W$.
The subfamily may have any size.  Each merger decreases the number of
blocks, so this procedure terminates after at most $m-1$ mergers.
Every resulting block satisfies
\begin{equation}\label{eq:block-terminal-properties}
 r(S)\le R_*,\qquad \ell(S)\le\zeta|S|,
\end{equation}
and no subfamily of at least two terminal blocks satisfies
\eqref{eq:block-merge}.

The anchor $i(S)$ need not attain the minimum defining $r(S)$.
Applying the triangle inequality through an index that does attain it gives
\begin{equation}\label{eq:block-anchor-radius}
 d(i,i(S))\le2\sqrt{r(S)}\le2\sqrt{\ell(S)/L}
 \qquad(i\in S).
\end{equation}
For singletons this holds with both sides zero.
By \eqref{eq:block-lens-input} and atomlessness, choose
$E_S\subseteq\bigcap_{i\in S}C_i$ of exactly the measure
\begin{equation}\label{eq:block-chosen-intensity}
 \lambda_S=e^{-\ell(S)}\bar c_{i(S)}.
\end{equation}
For singletons this choice is possible because $\bar c_i\le c_i$.
Thus $q_S=s_{i(S)}+\ell(S)\ge\log2$.
Different blocks have different anchor labels, and hence
\[
 \sum_Sq_S
 \le\sum_i s_i+\sum_S\zeta|S|
 =\sum_i s_i+\zeta m.
\]
This proves \eqref{eq:block-log-cost}.
For $n\ge3$, \eqref{eq:block-terminal-properties} and $m\le n^2$ give
\[
 \ell(S)\le L(1+R_*+\log m)\le L(B_0+3)\log n.
\]
Consequently \eqref{eq:block-anchor-loss} holds with $K=L(B_0+3)$.
Also
\[
 \log(1+\lambda_S^{-1})
 =q_S+\log(1+e^{-q_S})
 \le q_S+\log(3/2)
 \le\left(1+\frac{\log(3/2)}{\log2}\right)q_S,
\]
which proves \eqref{eq:block-first-budget}.
It remains to control the overlap rows.

\emph{What terminality implies.}
We claim that a constant $C_{\mathrm{ct}}=C_{\mathrm{ct}}(L,\zeta)$
bounds the number of blocks satisfying
\begin{equation}\label{eq:block-count-conditions}
 d(z,i(S))^2\le t,\qquad \ell(S)\le u
\end{equation}
by $C_{\mathrm{ct}}(1+t+u)$, provided that $z\in I$, $t,u\ge0$, and
\begin{equation}\label{eq:block-count-radius}
 8t+8u/L\le R_*.
\end{equation}
This is a count of blocks, including blocks with coinciding anchors in
pseudometric distance.

Let $h$ be the number of such blocks and let $W$ be their union, of size
$N$.  If $h\le1$ there is nothing to prove.
Otherwise choose one of their anchors, say $i(S_0)$.
For every $i\in S\subseteq W$, \eqref{eq:block-anchor-radius} gives
\[
 d(i(S_0),i)
 \le d(i(S_0),z)+d(z,i(S))+d(i(S),i)
 \le2\sqrt t+2\sqrt{u/L}.
\]
As $i(S_0)\in W$, it follows that
\[
 r(W)\le8t+8u/L=:D_*\le R_*.
\]
Choose a fixed integer $N_0\ge2$ so large that
$L\log N\le\zeta N/2$ for all $N\ge N_0$.
If
\[
 h\ge N_0\quad\text{and}\quad
 h\ge\frac{2L}{\zeta}(1+D_*),
\]
then $N\ge h$ implies
\[
 \ell(W)\le L(1+D_*+\log N)\le\zeta N.
\]
The whole subfamily could then be merged, contradicting terminality.
Since $1+D_*\le8(1+t+u)$, the claim follows, for example with
$C_{\mathrm{ct}}=1+N_0+16L/\zeta$.
Here the logarithm being absorbed is $\log N$, the logarithm of the
number of original labels in the union; the block sizes need not be
comparable.

\emph{Estimating a row.}
The elementary intersection bound and \eqref{eq:block-affinity-input}
give, respectively,
\begin{align}
 w_{ST}&\le e^{-|q_S-q_T|/2}\le1,
 \label{eq:block-weight-trivial}\\
 w_{ST}&\le C_0\exp\left\{
   \frac{\ell(S)+\ell(T)}2-c_0d(i(S),i(T))^2\right\}.
 \label{eq:block-weight-affinity}
\end{align}
For the first inequality, use
$\mu(E_S\cap E_T)\le\min(\lambda_S,\lambda_T)$ and discard the
factors $1+\lambda_S,1+\lambda_T$ in the denominator.
For the second, use $E_S\subseteq C_{i(S)}$ and
\eqref{eq:block-chosen-intensity}.

Fix $S$.  For $k\ge0$, put in the $k$th bin those $T\ne S$ for which
$e^{-(k+1)}<w_{ST}\le e^{-k}$; zero weights can be omitted.
Writing $Q=q_S+k+1\ge1$, the two bounds above imply
\begin{equation}\label{eq:block-bin-parameters}
 \ell(T)\le q_T<q_S+2(k+1)\le2Q,
 \qquad
 d(i(S),i(T))^2\le A_1Q,
\end{equation}
where $A_1=(\log C_0+2)/c_0$.
Indeed, taking logarithms in \eqref{eq:block-weight-affinity} gives
\[
 c_0d(i(S),i(T))^2
 <\log C_0+\frac{q_S+q_T}{2}+k+1
 \le\log C_0+q_S+2(k+1)
 \le(\log C_0+2)Q.
\]
Now set
\[
 A_2=\max\{1,8A_1+16/L\},\qquad \delta=A_2^{-1}.
\]
If $Q\le\delta R_*$, apply the terminality count with
$z=i(S)$, $t=A_1Q$, and $u=2Q$.
The number of blocks in the $k$th bin is at most
\begin{equation}\label{eq:block-bin-count}
 C_{\mathrm{ct}}(1+A_1Q+2Q)\le C_1(q_S+k+1),
\end{equation}
where $C_1=C_{\mathrm{ct}}(A_1+3)$.
These constants are independent of $B_0$.
Choose $B_0\ge\max\{1,6/\delta\}$, and increase $n_1$ so that
$R_*\ge4/\delta$ for $n\ge n_1$.

First suppose $q_S\le\delta R_*/4$, and put
$J=\lfloor\delta R_*/2\rfloor$.
For $0\le k\le J$ we have
$q_S+k+1\le\delta R_*$, so \eqref{eq:block-bin-count} bounds the
total weight in these bins by
\[
 C_1\sum_{k=0}^J e^{-k}(q_S+k+1)\le C_2(1+q_S).
\]
Every remaining weight is at most
$e^{-(J+1)}\le e^{-\delta R_*/2}$, and there are at most $n^2$ of them.
Their sum is therefore at most
\[
 n^2e^{-\delta R_*/2}=n^{2-\delta B_0/2}\le1.
\]
Since $q_S\ge\log2$, this proves
\begin{equation}\label{eq:block-row-cost}
 \log\left(1+\sum_{T\ne S}w_{ST}\right)\le C_3q_S
\end{equation}
in the present case.
If $q_S>\delta R_*/4$, the simpler bound
\eqref{eq:block-weight-trivial} gives
\[
 \log\left(1+\sum_{T\ne S}w_{ST}\right)
 \le\log(1+n^2)\le3\log n
 <\frac{12}{\delta B_0}q_S.
\]
Thus \eqref{eq:block-row-cost} holds for every block after enlarging $C_3$.
Summing it and applying \eqref{eq:block-log-cost} proves
\eqref{eq:block-row-budget}.
\end{proof}

Lemma~\ref{lem:caps} supplies the two hypotheses of
Proposition~\ref{prop:blocks} for the residual caps in the thin shell.
Its constants $C_0,c_0,L$ are independent of the fixed radius multiplier
$B_0$, so one first chooses these constants, then the value of $B_0$
specified in the proof, and finally the dimension threshold required by
Lemma~\ref{lem:caps}.
All constructions use the indexed list alone.  In particular, equal caps
remain separate labels, and for any fixed list the partition and the
measurable subsets may be chosen before the Poisson process is sampled.
